% Pitched slides (Oct 8 2026), kept outside the main deck. Build: from latex/, latexmk -pdf JMP_slides/pitched_slides_20261008.tex
% "What Moves Births" would sit after "A Loan Is Not Insurance"; the proposed Conclusion would replace the current one.
\input{JMP_slides/preamble.tex}
\title{Pitched Slides}
\begin{document}
% Pitched Oct 8: summary of the fixed-price mechanism results above (numbers from the two preceding frames, base 14.402).
\begin{frame}{What Moves Births}
\small
Same households and parameters, house prices fixed:
\begin{itemize}\setlength\itemsep{0.7em}
  \item \textbf{The cost of space.} House prices and rents $+10\%$: completed fertility $-5.3\%$.
  \item \textbf{Income risk.} Lower risk: first-birth attempts $+39\%$, childlessness at 45 from 19\% to 7\%. The effect grows with the cost of space ($+33\%$, $+39\%$, $+44\%$).
  \item \textbf{Credit.} Financed share $80\%\to95\%$: ownership at 18--29 up 25 points, completed fertility $-1.4\%$. Renter credit moves births earlier, not more.
  \item \textbf{Insurance.} A transfer paid in bad years raises completed fertility $14.7\%$; the same outlay paid to everyone, $6.6\%$.
\end{itemize}
\end{frame}

\begin{frame}{Conclusion (proposed)}
% Proposed rewrite; the author's original bullets are kept as a comment below.
%  \item A lifecycle model of fertility with realistic housing: tenure, down payments, space
%  \item Children need space; large homes must be bought; the old hold most of them
%  \item Today is a point on a transition to a smaller population
%  \item A rebated property tax can move space toward young families, generating a fertility response
\begin{itemize}\setlength\itemsep{1.0em}
  \item A lifecycle model of fertility with realistic housing: children need space, large homes must be bought, buying needs a down payment
  \item The cost of family space matters: 10\% dearer space lowers completed fertility by about 5\%
  \item Income risk is first order: a child commits the household to space, and lower risk raises first-birth attempts by about 40\%
  \item A loan is not insurance: more mortgage credit raises young ownership, not family size; a transfer paid in bad years does twice as much as the same outlay paid to all
  \item A rebated property tax lowers house prices in the long run and supports a slightly larger population
\end{itemize}
\end{frame}
\end{document}
